\chapter{Ground: Outlines and Their Construction}
\label{ch:ground}

\section{Islands and the Height Convention}

The ground level deals in \textbf{shapes} on \textbf{layers}, grouped into
\textbf{groups}. A shape has an outline, a floor and a thickness, and where its
ground touches or overlaps another shape's in the same group, the two build one
\textbf{island}.

The first thing to get right is the height convention, because it is off by one
from the obvious reading and an agent gets it wrong every time otherwise:
\id{base\_height: 9} puts the top block at $y=8$, not $y=9$. A column read at such
a shape answers grass at $y8$, dirt at $y7$ and $y6$, stone from $y5$ to $y1$, and
bedrock at $y0$.

\section{Four Forms of Outline}

\subsection{Polygons and the Rectangle}

A polygon is a ring of points. A rectangle is not: it states bounds rather than
points, and so it is restated as a four-point polygon before any point of it can
move. That restatement is invisible in the document and load-bearing in practice.

Three routes edit one point without disturbing any other:

\begin{lstlisting}
PATCH  /api/map/{slug}/sketch/shapes/{id}/vertices/{index}   {x, z}
POST   /api/map/{slug}/sketch/shapes/{id}/vertices           {after, x, z}
DELETE /api/map/{slug}/sketch/shapes/{id}/vertices/{index}
\end{lstlisting}

The insert has index arithmetic that bites. A point added \id{after: 0} lands at
index 1, and every later index goes up by one, so a pair of points previously at
2 and 3 migrates to 3 and 4. An agent replaying a list of edits by index, computed
against the original ring, will scramble the outline.

\subsection{Height as a Per-Vertex Field}

\id{anchor\_heights} is one thickness per vertex, in the same ring order, and the
surface between them is interpolated across the face. A six-point island carrying
\id{[9, 12, 16, 17, 17, 14]} is the plane $9 + x/4 + z/4$ in the island's own
frame. That is a tilt, authored with six numbers.

\subsection{Curves and B\'ezier Handles}

Four ways to state an outline that is not made of straight edges. A \textbf{circle}
takes a centre and a radius and the studio draws it as a 64-sided polygon. A
\textbf{lasso} is, in the document, simply a polygon. The difference is the
canvas, which traces it freehand and simplifies the trace at a four-block
tolerance. And any polygon edge can be bowed with Bézier handles.

The Bézier contract is worth stating exactly, because every part of it is a place
to go wrong:

\begin{quote}
\id{controls} is keyed by a vertex's index as a string, and every handle is an
\textbf{absolute point on the board}. The edge from vertex $i$ to vertex $i+1$ is a
cubic curve from vertex $i$, pulled by \id{controls[i].out} and
\id{controls[i+1].in}, to vertex $i+1$. So a vertex's \id{out} bends the edge after
it, and its \id{in} bends the edge before it. The studio samples every curved edge
16 times, and the curve still passes through every vertex.
\end{quote}

Bowing one edge is two handles:

\begin{lstlisting}[language=json]
"controls": { "0": { "out": [74.67, -6] }, "1": { "in": [81.33, -6] } }
\end{lstlisting}

Rounding a whole ring is a closed-form recipe rather than a hand fit: at each
vertex $P$ with neighbours $\mathit{previous}$ and $\mathit{next}$,
\[
\mathit{in} = P - \frac{\mathit{next} - \mathit{previous}}{6}, \qquad
\mathit{out} = P + \frac{\mathit{next} - \mathit{previous}}{6}.
\]

\begin{ruling}{Limits on a B\'ezier Handle}
Keep each handle at least as far \emph{along} its edge as it is \emph{away} from
it, and keep the bulge under about a third of the edge's length. Past that the
curve crosses itself.
\end{ruling}

\subsection{Polylines}

A polyline is a line with a width: its points are a centreline, and the band
\id{radius} blocks to each side of it is ground. The order of operations is the
fact to know:

\begin{quote}
The points are \textbf{smoothed first}, by a centripetal Catmull--Rom spline with
eight samples between each pair, and only then offset \id{radius} blocks to either
side, so four points draw a curve rather than three chords. The ends are cut
square.
\end{quote}

\begin{lstlisting}[language=json]
{"id":"line-solid","type":"polyline","operation":"add","floor":0,
 "base_height":9,"vertices":[[0,10],[13,3],[27,17],[40,10]],
 "radius":3,"stroke_edge":"solid","stroke_seed":7}
\end{lstlisting}

\id{stroke\_edge} takes three values with numbers attached. \id{solid} keeps one
width. \id{rough} lets it swing up to 45\% either way with the two sides wandering
apart, seeded by \id{stroke\_seed}. \id{tapered} narrows it to 35\% at the ends.
A polyline also takes \id{anchor\_heights}, which are heights at its points, with
every block taking the height of the nearest place along the line. A lane can
therefore rise.

Two facts separate a polyline from things that look like it. A polyline whose ends
lie \emph{inside} two islands joins them into one, which is how a land bridge is
drawn. And, in the distinction that costs an afternoon, a polyline shape \emph{is
ground}, whereas the \id{stroke} prop, which looks alike, repaints the top course
of what it crosses and adds no ground at all.

\begin{ruling}{Self-Overlap in a Polyline}
The band is one outline filled even--odd, so where its two sides cross they cancel
and build void. The studio names it \id{SK25}. Four points across 26 blocks at
radius 3 lap in every bend; the same four points spread over 40 blocks clear. A
coil is drawn as one polyline per part-turn.
\end{ruling}

\section{Overlap: Height and Paint}

Where two shapes cover a cell, two questions are settled separately, and conflating
them is a common error.

\begin{quote}
The height and the paint of an overlap are decided separately. Where two shapes
cover a cell, the taller one gives the column its height, and the paint goes to the
\textbf{smallest themed shape whose top is that tallest top}.
\end{quote}

So on a flat island where a yard and a meadow both stand at 9, the yard being
smaller means the overlap paints as the yard. A lawn smaller again paints as the
lawn inside the yard. Height is decided by tallest; paint, among the shapes that
tie for tallest, by smallest.

\begin{ruling}{Occlusion Within a Layer}
On one layer the taller shape takes the whole column, so a lower shape drawn inside
a higher one simply is not there, and the studio names it \id{SK9}. A sunken place
inside a raised one is drawn as shapes \emph{around} it, or on a layer of its own. This is the same rule \cref{ch:layers} states from the other side.
\end{ruling}

\section{Faces and Band Stacks}

Where a shape's ground drops, whether over lower ground or over the void, the
exposed side is a \textbf{face}, and a theme can stripe it with a band stack.

One property of that stack is not what it looks like:

\begin{quote}
A band stack does not cycle: past its last band, \id{"ending": "repeat"} hands
every further course to that last band.
\end{quote}

So striping a face sixteen courses tall with a four-course pattern means writing
the four courses out four times. A stack that reads as four bands in the author's
head is twelve bands in the document.

The other property is that only the courses standing \emph{above} the neighbouring
ground are face. The same yard, stepped up from a meadow on one side and dropping
to the void on the other, reads four striped courses on the first and the whole
stripe from top to bedrock on the second.

\section{Five Forms of Ascent}

Getting from one height to another is its own small craft, and the studio offers
five forms: a ramp (one polygon with anchors at each end), a stair (one polygon
with deeper treads), a turned flight (two polygons with a landing between), a run
of steps (one rectangle per tread), and a floating flight (one raised slab per
tread, with air under each).

One rule governs all of them and it is the most counter-intuitive fact in the whole
ground model:

\begin{quote}
\id{anchor\_heights} are heights at a polygon's corners, and every block reads the
slope at its own centre and rounds. A flight therefore comes out even \textbf{only
when the anchors sit half a riser past its first and last tread}: a 2-deep stair
whose treads are 1 to 7 states 0.5 and 7.5. Whole-number anchors 8 and 1 over 16
blocks build treads 1, 2 and 3 deep. The ramp works with whole numbers because one
block up per block rounds cleanly.
\end{quote}

And a consequence that decides how a flight can be authored at all: the studio's
canvas rounds anchor heights to whole blocks, so the half-block form can only be
written into the document. Drawing an even stair by hand means one rectangle per
tread instead.
